\begin{abstract}
“八百里巢湖，烟波浩渺”，巢湖沿岸湿地植被如绿绸环绕，既是水陆交锋的“绿色卫士”，亦是区域生态平衡的命脉所系。本文以环巢湖湿地植被为研究对象，系统剖析其群落结构与环带状分布格局——从沉水植物的“水下森林”到挺水植物的“青纱帐”，从浮叶植物的水面遮蔽到湿生草甸的滩涂固持，植被分布深刻响应水文梯度、土壤特性与人类活动的耦合驱动。

在此基础上，进一步阐释湿地植被的多维生态福祉：净水除污以控富营养化，生灵栖居以维生物多样性，固碳释氧以助气候调节，护岸蕴文以兼生态与文化价值。然而，当前湿地植被面临水利扰序、物种入侵、污染承压、生境破碎的多重威胁。为此，本文提出植被重建、水文调控、污染拦截、社区共治的综合修复路径，旨在为守护巢湖生态底色、实现人与自然和谐共生提供科学支撑，让环巢湖湿地永葆“渚清沙白鸟飞回，水草丰美鱼虾肥”的盛景\cite{chu2023,tang2022,liu2022}。

\textbf{关键词：环巢湖湿地；植被群落；生态功能；保护修复；可持续发展}
\end{abstract}
